\begin{figure}[!hbp]
\centering
\begin{tikzpicture}[x=.48cm,y=.48cm,font=\small,
 ordinary/.style={draw=ink,line width=.7pt},
 crossing/.style={draw=rust,densely dashed,line width=1.15pt},
 chosen/.style={circle,draw=ink,fill=ink,inner sep=0pt,minimum size=4.3pt},
 outside/.style={circle,draw=ink,fill=white,inner sep=0pt,minimum size=4.3pt}]
\begin{scope}
  \foreach \i in {0,...,3} {
    \draw[ordinary] (\i,0)--(\i,3);
    \draw[ordinary] (0,\i)--(1,\i) (2,\i)--(3,\i);
    \draw[crossing] (1,\i)--(2,\i);
  }
  \foreach \i in {0,...,3} {
    \foreach \j in {0,1} \node[chosen] at (\j,\i) {};
    \foreach \j in {2,3} \node[outside] at (\j,\i) {};
  }
  \node at (1.5,-1.05) {Square grid};
\end{scope}
\begin{scope}[xshift=3.65cm]
  \coordinate (root) at (1.8,3);
  \coordinate (l) at (.6,1.5); \coordinate (r) at (3,1.5);
  \coordinate (ll) at (0,0); \coordinate (lr) at (1.2,0);
  \coordinate (rl) at (2.4,0); \coordinate (rr) at (3.6,0);
  \draw[crossing] (root)--(l);
  \draw[ordinary] (root)--(r) (l)--(ll) (l)--(lr) (r)--(rl) (r)--(rr);
  \foreach \v in {l,ll,lr} \node[chosen] at (\v) {};
  \foreach \v in {root,r,rl,rr} \node[outside] at (\v) {};
  \node at (1.8,-1.05) {Complete binary tree};
\end{scope}
\begin{scope}[xshift=7.35cm]
  \draw[ordinary] (0,1.5)--(2.25,1.5) (3,1.5)--(4.5,1.5);
  \draw[crossing] (2.25,1.5)--(3,1.5);
  \foreach \i in {0,...,3} \node[chosen] at (.75*\i,1.5) {};
  \foreach \i in {4,...,6} \node[outside] at (.75*\i,1.5) {};
  \node at (2.25,-1.05) {Path};
\end{scope}
\end{tikzpicture}
\caption{Three cuts between filled and open vertices. Four edges
cross the middle of this grid. On the binary tree a balanced split
can cross just one edge, so the ABP argument prescribes the size of
the smaller set more carefully. On the path, a proper nonempty prefix
of any size cuts just one edge. The formula argument combines splits
at several scales.}
\label{fig:witness-cuts}
\end{figure}
